% Section III -- Methodology.
% Numbers here come from `python analyze.py results/ --paper` on the 18-run
% sweep (6 epsilon x 3 seeds). If a run is added or a control changes,
% regenerate that output and update this file rather than editing by hand.
% One figure is left to verify against this sweep -- see the VERIFY comment
% in the Energy Measurement subsection.
% Citation keys are placeholders: \cite{opacus}, \cite{flower}, \cite{rdp},
% \cite{dpsgd}, \cite{dirichlet}, \cite{nvml}.

\section{Methodology}

\subsection{Federated Setup}

We train a five-layer convolutional network (0.24~MB) on CIFAR-10 across 20
simulated clients using the Flower framework~\cite{flower}. Each round samples
50\% of clients for training and 50\% for evaluation. Client data is
partitioned by a Dirichlet distribution over labels with $\alpha=0.5$, the
standard non-IID setting~\cite{dirichlet}; partition sizes therefore differ
between clients, which matters for privacy accounting below. Clients perform
one local epoch of SGD per round (batch size 32, learning rate 0.01, momentum
0.9) and the server aggregates by FedAvg over 30 rounds. Every configuration
is repeated with three seeds.

\subsection{Privacy Accounting}

We enforce \emph{client-level local differential privacy}: each client holds
its own $(\varepsilon,\delta)$ budget over the entire federated run, and the
server is untrusted.

The budget is accounted over all local SGD steps a client takes across all
rounds in which it participates --- not per round, and not per local epoch.
For a client with $n$ local training examples, batch size $B$, $R$ rounds,
client sampling fraction $q_c$ and $E$ local epochs:
%
\begin{align}
  s_{\text{epoch}} &= \lceil n / B \rceil \\
  s_{\text{total}} &= R \, q_c \, E \, s_{\text{epoch}}
\end{align}
%
We invert the RDP accountant~\cite{rdp} once, before training starts, to find
the noise multiplier $\sigma$ that yields the target $\varepsilon$ at $\delta$
after $s_{\text{total}}$ steps, and hold $\sigma$ constant for the run.

Precomputation is deliberate. In simulation each client object is
reconstructed every round, so an engine-internal accountant resets each round
and reports a per-round $\varepsilon$ while the true budget compounds ---
understating the privacy loss by roughly a factor of the round count.

Two properties of the sampling deserve statement. First, Opacus'
\texttt{DPDataLoader} samples at rate $1/\lceil n/B \rceil$ rather than $B/n$;
we calibrate at the rate the sampler actually uses, so the realised
$\varepsilon$ matches the target to within the accountant's search tolerance
rather than being conservative by an unstated margin. Second, we claim
\textbf{no privacy amplification by client subsampling}: amplification
requires secure aggregation or a shuffler, neither of which we assume. The
reported $\varepsilon$ is therefore an upper bound.

Because Dirichlet partitions differ in size, $\sigma$ differs per client.
Table~\ref{tab:sigma} reports the observed range.

\begin{table}[t]
\caption{Per-client noise multiplier, 30 rounds, $\delta=10^{-5}$}
\label{tab:sigma}
\centering
\begin{tabular}{lcc}
\hline
$\varepsilon$ & $\sigma_{\min}$ & $\sigma_{\max}$ \\
\hline
0.5 & 2.773 & 6.094 \\
1   & 1.592 & 3.320 \\
2   & 1.038 & 1.899 \\
4   & 0.781 & 1.204 \\
8   & 0.621 & 0.853 \\
\hline
\end{tabular}
\end{table}

\subsection{Energy Measurement}

Energy is read from the NVML hardware counter
\texttt{nvmlDeviceGetTotalEnergyConsumption}~\cite{nvml}, which accumulates
millijoules on the device. We difference the counter before and after each
round, so the measurement involves no power sampling and no numerical
integration, and cannot miss a transient. This is a measured quantity, not a
TDP-based estimate.

Measurement is server-side and per round, bracketing \texttt{configure\_train}
through \texttt{aggregate\_evaluate}. In simulation all clients share one GPU,
so a round-level window captures the full cost regardless of how the runtime
schedules client actors; measuring inside each client would double-count
concurrent work.

We report energy both gross and net of an idle baseline. The baseline is
sampled over a quiet-GPU window before each run: $18.54$~W (sd $0.10$) on our
hardware. Net energy subtracts this baseline over the wall-clock the counter
actually covers, which is \emph{not} whole-process wall-clock --- run
start-up, centralised evaluation between rounds, and teardown fall outside the
round windows. Subtracting idle over the longer interval would understate net
compute energy substantially.% VERIFY the exact fraction on this sweep:
% python -c "import json,glob;rs=[json.load(open(p)) for p in glob.glob('results/run_*.json')];
%   print(sum(1-r['totals']['wall_s_measured']/r['totals']['wall_s_total'] for r in rs)/len(rs))"
% then state it: "...and account for roughly N% of process wall-clock."

Repeating each configuration across three seeds gives a measurement noise
floor of $\pm 4.38$~W ($6.38\%$). We report no effect smaller than this as an
effect.

\subsection{Experimental Controls}

\textbf{Fixed clipping norm.} $C=1.0$ for every $\varepsilon$. Tuning $C$ per
budget would confound the privacy--utility comparison. Our ablation shows
$C$ affects convergence substantially while leaving per-round energy
unchanged, so $C=1.0$ is a conservative rather than optimal choice.

\textbf{Identical compute across budgets.} The energy axis is only
interpretable if every $\varepsilon$ condition executes the same workload. We
verify per round that executed optimiser steps match the steps the accountant
was calibrated for. Opacus' loader yields exactly $\lceil n/B \rceil$ batches
per epoch, so step counts are deterministic and independent of
$\varepsilon$; Poisson sampling randomises batch \emph{size} (approximately
$\pm 2\%$ of $n$), not step count, and that residual is unbiased with respect
to $\varepsilon$. Across the sweep, executed and calibrated steps agree to
within one step per round ($0.16\%$), a floating-point artefact of the
loader's batch-count computation that runs \emph{fewer} steps than accounted,
preserving the upper-bound claim.

\textbf{Accountant excluded from the measurement window.} Deriving $\sigma$
requires a binary search over the RDP accountant whose cost depends strongly
on the target budget: we measure $0.086$~s at $\varepsilon=0.5$ against
$0.693$~s at $\varepsilon=8$. Executed inside the federated loop, this cost
falls within the energy-measurement window and scales with $\varepsilon$,
producing a wall-clock trend that is indistinguishable in shape from a cost of
privacy. In our first clean sweep it accounted for the entire measured
$\varepsilon$ trend (predicted $+18.2$~s at $\varepsilon=8$ relative to
$\varepsilon=0.5$; measured $+19.7$~s) while step counts stayed flat. We
therefore derive all $\sigma$ values before measurement begins. We flag this
as a general hazard for energy studies of privacy mechanisms.

\textbf{Randomised run order.} Sweeping $\varepsilon$ in a fixed order makes
execution order and $\varepsilon$ rank-correlated ($\rho=-0.99$), so any drift
over a sweep is inseparable from a privacy effect. We randomise the order of
all 18 runs, reducing the correlation to $\rho=-0.20$ and making the
$\varepsilon$ trend attributable.

\textbf{Determinism deliberately not enforced.} We do not enable
\texttt{torch.use\_deterministic\_algorithms}: it forces slower kernel
selection and would change measured energy, i.e.\ we would be measuring a
configuration nobody deploys. Reproducibility comes from averaging over seeds.
Per-client step counts are deterministic; the round-to-round variation in
aggregate step count comes from client sampling, which we record per round.

\subsection{Limitations}

Clients are virtual and share a single GPU rather than running on distinct
devices. Relative comparisons between privacy budgets hold; absolute
per-device joules do not transfer to real edge hardware. CPU RAPL counters
were unavailable in our containerised environment, so measurement is GPU-only
and communication energy is modelled analytically rather than measured. The
workload does not saturate the device (roughly 68~W on a 360~W card), which is
why idle-subtracted figures are reported alongside gross ones.
